\documentclass[border=4pt]{standalone}
\usepackage{amsmath,amssymb,mathtools,xcolor,varwidth}
\definecolor{resultblue}{HTML}{1E6FE0}
\definecolor{radgreen}{HTML}{00A35A}
\definecolor{ellblue}{HTML}{1E6FE0}
\begin{document}
\begin{varwidth}{28em}
\large\bfseries
Compounding the pieces, the rectangle $PvG : Qv^2 = PC^2 : CD^2$, and by the      similar triangles $QxT,\,PEF$ we get $QT^2 : EP^2 = PF^2 : CA^2$; joining these      with Lemma XII forces \textcolor{resultblue}{$L\cdot QR / QT^2 = SP^2$}. Substituting into      the Prop. VI measure $F \propto 1/(SP^2 \cdot QT^2/QR)$ makes it become      $1/(SP^2\cdot SP^2)$ divided cleanly, so the \textcolor{resultblue}{force varies as      $1/SP^2$}. Thus the \textcolor{radgreen}{distance $SP$} to the focus alone governs the pull:      the same \textcolor{resultblue}{inverse-square law} that Newton would recognise as universal      gravitation holds every planet in its \textcolor{ellblue}{Keplerian ellipse}. \textit{Q.E.I.}
\end{varwidth}
\end{document}
